\begin{abstract}
\noindent The leave-out estimator of Kline, Saggio and S{\o}lvsten (2020, KSS) removes the bias of variance
components in two-way fixed effects wage models. Its sampling variance depends on the error variance of every
observation, which a single observation cannot reveal. This note proposes a standard error built on one rule:
estimate each error variance from independent repeat measurements of the same mean wherever they exist, using
only the independence that the point estimator already assumes, and use a model elsewhere. The wages of a
\emph{stayer} in different years are such repeat measurements, so every stayer gets its own variance estimate,
with no model linking it to other workers. \emph{Movers} get a model: their error variances are a scale, which depends on the length of the match and
on the leverage, times a quadratic function of the estimated person and firm effects, both fitted to the
movers' own leave-out variance estimates. We show that the variance estimator is exactly unbiased when movers'
variances are known, however variances differ across stayers; that it converges under moment bounds and under
design conditions implied by those KSS use for asymptotic normality; and that with the model for movers,
evaluated at known characteristics of workers and firms, it converges to the truth plus an explicit term,
which is zero when the model holds. The implementation evaluates the model at estimated person and firm
effects; for that step we rely on simulations. The computations reuse what FEVC already computes for the point
estimate and add a small number of linear solves. In simulations on the design of a six-year Veneto panel
(\R{real.n.rows} rows), with estimated effects throughout, the root-mean-square standard error is between
\R{sim.all.min} and \R{sim.all.max} times the true one in ten scenarios, and the coverage of the 95\% interval
is between \R{sim.all.cov.min} and \R{sim.all.cov.max}. In these runs the standard error is more often too large than too small.
This is the net of opposing contributions, which the note measures, and not a proven property. A
validation with tolerances fixed in advance, on fresh draws and on two smaller designs cut from the same
data, met them on two designs; on the third, a three-year window, the standard error of the variance of firm
effects was up to \R{val.all.over.max} too large, beyond the tolerance of 15\%. Coverage in these runs is
between \R{val.all.cover.min} and \R{val.all.cover.max}. For
designs with one weakly identified direction, the note shows that the same variance estimates, with KSS's own
estimate for the variance of that direction, give a consistent covariance matrix for the interval of KSS, and
that a rule which guards that estimate against being too small keeps the interval's coverage bound. In
simulations on synthetic panels in which two groups of firms are linked by 10 to 100 workers, or not divided
at all, that interval covers between \R{q1.acc.min} and \R{q1.acc.max}. The normal interval covers as little as
\R{q1.normal.worst} there, in the worst cases because its variance estimate is not positive in up to
\R{q1.vneg.max} of the draws. One of the tolerances set in advance, on the length of the interval without a
bottleneck, was missed.
\end{abstract}
